\chapter{G 结构的几何学}

% [原文 p.36]
在 Erlangen 大学的教授就任演讲中，年轻的 Felix Klein 提出了如下观点：对空间 $M$ 与其上可迁地作用的群 $G$ 构成的组 $(M,G)$，对应着一种几何学，只有在该群作用下不变的空间性质才成为几何学的研究对象。例如，Euclid 几何对应 Euclid 空间与 Euclid 合同群，仿射几何对应仿射空间与仿射变换群，而射影几何对应射影空间与射影变换群。仿射变换群是由保持无穷远超平面的元素构成的射影变换群的子群，Euclid 变换群是由保持 Euclid 距离的元素构成的仿射变换群的子群。在较大群下保持不变的性质自动也在较小群下保持不变，反之则未必成立，因此对应于较大群的几何学被认为是更基本的。Klein 由此把当时已知的各种几何学安置到以射影几何为顶点的层级之中，同时也促进了 Lie 群与齐性空间的研究。
% [Tutor Note: 「Erlangen 大学的教授就任演讲」即 Klein 1872 年的埃尔朗根纲领（Erlanger Programm），其中提出以变换群对几何学分类的思想。]

把 Klein 流派的这一观点扩展到包含 Riemann 几何的形式，就是 Élie Cartan 的 G 结构（G-structure）几何学。对 Euclid 空间 $\R^n$ 与其上作用的 Lie 群 $G$ 构成的组 $(G,\R^n)$，$n$ 维流形 $M$ 的 G 结构是指 $M$ 的切丛 $TM$ 的结构群 $\GL(n,\R)$ 到 $G$ 的约化。也就是说，取 $M$ 的开覆盖，把切丛 $TM$ 表示为平凡丛的粘合时，转换函数的值取在 $G$ 中而非 $\GL(n,\R)$ 中。当 $G$ 如旋量群那样不是 $\GL(n,\R)$ 的子群时，与其说是「约化」，毋宁说有「提升」的意味。代表性的 G 结构有 $G=\GLp(n,\R),\,\OO(n),\,\GL(n/2,\C)$ 等情形，分别对应可定向流形、Riemann 流形、近复流形。此外，G 结构常由张量场刻画，上述三个例子中分别对应体积形式（处处不为零的最高次微分形式）、Riemann 度量（处处正定的 2 阶对称张量）、近复结构（平方为 $-1$ 倍的切丛自同构）。

% [原文 p.37]
G 结构中重要的概念是可积性。流形由 Euclid 空间的开集粘合而成，粘合映射的微分给出切空间之间的线性映射；当能取到使该映射总在 $G$ 中的图册时，称此流形具有可积的 G 结构。换一种说法：$\R^n$ 具有标准的 G 结构；具有 G 结构的流形，其各点的切空间与 $\R^n$ 连同 G 结构一起被等同；而具有可积 G 结构的流形，其各点处有限大小的邻域与 $\R^n$ 的开集连同 G 结构一起被等同。例如，可积的 Riemann 流形由 $\R^n$ 的开集通过等距映射粘合得到，这样的流形称为平坦流形；由 Bieberbach 定理，已知它若紧致，则只能是环面的有限群商。另一方面，可积的近复流形称为复流形，由 $\C^n$ 的开集通过双全纯映射粘合得到。双全纯映射与微分同胚相比是「硬的」，但比等距映射「软」，因此复流形的数量远多于平坦流形。G 结构的可积性也可以用刻画 G 结构的张量的语言来叙述。例如，Riemann 流形平坦的充要条件是 Riemann 曲率张量为零，而近复结构可积的充要条件是 Nijenhuis 张量为零。
% [Tutor Note: 由 Bieberbach 定理（1911–1912），紧致平坦流形必为环面被有限群自由作用所得之商。]

本章从 G 结构几何学的观点概观辛几何及与之相关的几何。

\section{辛向量空间}

对自然数 $n$，用 $I_n$ 表示 $n$ 阶单位矩阵，定义 $2n\times 2n$ 矩阵 $J_{2n}$ 为
\begin{equation}
  J_{2n}=\begin{pmatrix}0 & I_n\\ -I_n & 0\end{pmatrix}.
  \tag{3.1.1}
\end{equation}

\begin{proposition*}{3.1.1}
  对实向量空间 $V$ 上任意的交错双线性形式 $\omega$，存在自然数 $r$ 与 $V$ 的适当基底，使得在此基底下 $\omega$ 可表示为分块对角矩阵
  \begin{equation}
    \begin{pmatrix} J_{2r} & \\ & 0\end{pmatrix}
    \tag{3.1.2}
  \end{equation}
  的形式。
\end{proposition*}

\begin{proof}
  使用 Gram--Schmidt 正交化法（Gram-Schmidt orthogonalization）。首先把 $\omega$ 视为映射 $V\to V^*$，取其核 $\ker\omega$ 的一个基底 $\{g_i\}_{i=1}^{k}$ 并固定。接着取 $\ker\omega$ 的一个补空间，记为 $V_1$，则 $\omega$ 在 $V_1$ 上确定一个交错非退化双线性形式。这里取 $V_1$ 的非零元 $e_1$，由于 $\omega$ 在 $V_1$ 上非退化，存在某个 $f_1\in V_1$ 使 $\omega(e_1,f_1)=1$。把 $\{e_1,f_1\}$ 生成的 $V_1$ 的子空间关于 $\omega$ 的正交补空间记为 $V_2$。从 $V_1$ 到 $V_2$ 的正交投影为
  \begin{equation}
    v\mapsto v-\omega(v,f_1)e_1+\omega(v,e_1)f_1
    \tag{3.1.3}
  \end{equation}
  % [原文 p.38]
  所给出。由于 $\omega$ 在 $V_2$ 上的限制仍非退化，若 $V_2\neq 0$，则存在 $\{e_2,f_2\}\subset V_2$ 使 $\omega(e_2,f_2)=1$。反复进行此操作，得到 $V$ 的基底 $\{e_1,\dots,e_\ell,f_1,\dots,f_\ell,g_1,\dots,g_k\}$。由构造，$\omega$ 在此基底下的表示由 (3.1.2) 给出。
\end{proof}

\begin{corollary}{3.1.2}
  向量空间上交错 2 次形式的秩为偶数。
\end{corollary}

实向量空间 $V$ 与其上非退化的 2 次形式 $\omega$ 的组 $(V,\omega)$ 称为辛向量空间（symplectic vector space）。$V$ 的维数必为偶数，并且 $\omega$ 在 $V$ 的适当基底 $\{e_1,\dots,e_n,f_1,\dots,f_n\}$ 下由矩阵 $J_{2n}$ 表示：
\begin{equation}
  \omega(e_i,e_j)=\omega(f_i,f_j)=0,\quad \omega(e_i,f_j)=\delta_{ij},\quad i,j=1,\dots,n.
  \tag{3.1.4}
\end{equation}
这样的基底称为 $(V,\omega)$ 的辛基底（symplectic basis）。

对辛向量空间 $(V,\omega)$，由保持 $\omega$ 的 $V$ 的线性自同构组成的 $\GL(V)$ 的子群称为辛群（symplectic group），记作 $\Sp(V,\omega)$。使用 $(V,\omega)$ 的辛基底，辛群可表示为
\begin{equation}
  \Sp(V,\omega)=\{g\in\GL(2n,\R)\mid g^{\mathsf T}J_{2n}g=J_{2n}\}.
  \tag{3.1.5}
\end{equation}

\begin{remark}{3.1.3}
  辛群（symplectic group）最初由 Weyl 命名为“complex group”，但因与「定义在复数域上的群」相混淆，于是把源自拉丁语的 complex 一词的构成要素 ‘com’ 与 ‘plexus’ 分别替换为希腊语中对应的构成要素 ‘sym’ 与 ‘plectikos’，Weyl 由此新造了 symplectic 一词。
  % [Tutor Note: Weyl 于 1939 年的《The classical groups》中自述此事；symplectic 与 complex 同源，本义为「编织在一起的」。]
\end{remark}

\takeaway{辛向量空间是带有非退化交错 2 次形式的偶维向量空间。}

设 $(V,\omega)$ 为辛向量空间。对 $V$ 的子空间 $W$，其正交补空间（orthogonal complement）为
\begin{equation}
  W^\perp=\{v\in V\mid \text{对任意 } w\in W \text{ 都有 } \omega(v,w)=0\}
  \tag{3.1.6}
\end{equation}
$W$ 称为：
\begin{itemize}[leftmargin=2.5em,itemsep=2pt,topsep=3pt]
  \item 满足 $W\subset W^\perp$ 时的各向同性子空间（isotropic subspace），
  \item 满足 $W\supset W^\perp$ 时的包合子空间（involutive subspace），
  \item 满足 $W=W^\perp$ 时的 Lagrange 子空间（Lagrangian subspace）。
\end{itemize}
% [Tutor Note: 原书用「包合的（involutive）」，即通称的余各向同性子空间（coisotropic subspace）；此处照原文译名。]

% [原文 p.39]
\begin{proposition*}{3.1.4}
  $(V,\omega)$ 为辛向量空间，$W\subset V$ 为各向同性子空间时，
  \begin{equation}
    \dim W\le \frac{1}{2}\dim V
    \tag{3.1.7}
  \end{equation}
  成立。
\end{proposition*}

\begin{proof}
  使用命题 3.1.1 的证明中所用的 Gram--Schmidt 正交化法，对 $W$ 的任意基底 $\{e_1,\dots,e_k\}$，可以取到包含它的 $V$ 的辛基底 $\{e_1,\dots,e_n,f_1,\dots,f_n\}$，因此
  \begin{equation}
    2\dim W=2k\le 2n=\dim V
    \tag{3.1.8}
  \end{equation}
  成立。
\end{proof}

\section{Darboux 定理}

赋予 $2n$ 维可微流形 $M$ 以 Cartan 意义下的 $\Sp(2n,\R)$ 结构与赋予 $M$ 上一个处处非退化的 2 次微分形式 $\omega$ 是等价的。下面的定理是关于辛流形局部结构的基本定理：
% [Tutor Note: 「Cartan 意义下的 G 结构」指切丛的结构群约化到 G 的结构；这里 Sp(2n,R) 结构即辛结构。]

\begin{theorem}{3.2.1（Darboux）}
  对任意流形 $M$ 上任意的非退化 2 次微分形式 $\omega$，由 $\omega$ 决定的 $\Sp(2n,\R)$ 结构可积与 $\omega$ 为闭形式是等价的。
\end{theorem}

因此，辛几何是研究可积的 $\Sp(2n,\R)$ 结构的学问。或者按 Klein 的风格，也可以说是研究辛流形的辛同胚不变性质的学问。

Darboux 定理常不用定理 3.2.1 的形式，而用与之等价的下列形式叙述：

\begin{theorem*}{3.2.2（Darboux）}
  对辛流形 $(M,\omega)$ 上的任意点 $p$，存在 $p$ 的邻域 $U$、$\R^{2n}$ 的开集 $V$，以及微分同胚 $\psi\colon V\to U$，使得
  \begin{equation}
    \psi^*\omega=\sum_{i=1}^{n}\rd p_i\wedge \rd q_i
    \tag{3.2.1}
  \end{equation}
  成立。
\end{theorem*}

为证明 Darboux 定理，使用下述定理：

\begin{theorem*}{3.2.3（Moser）}
  设 $M$ 为紧致流形，$(\omega_t)_{t\in[0,1]}$ 为以闭区间 $[0,1]$ 参数化的 $M$ 上辛形式的光滑族。此时，若 $\omega_t$ 所属的 de Rham 上同调类不依赖于 $t$，则存在 $M$ 的
  % [原文 p.40]
  微分同胚的光滑族 $(\varphi_t)_{t\in[0,1]}$，满足 $\varphi_0=\id_M$，且对任意 $t\in[0,1]$ 有
  \begin{equation}
    \varphi_t^*\omega_t=\omega_0.
    \tag{3.2.2}
  \end{equation}
\end{theorem*}

\begin{proof}
  给定满足 $\varphi_0=\id_M$ 的 $M$ 的微分同胚族 $(\varphi_t)_{t\in[0,1]}$ 时，对任意 $t\in[0,1]$ 使 (3.2.2) 成立的充要条件是：把 (3.2.2) 的左端对 $t$ 微分所得的
  \begin{align}
    \frac{d}{dt}(\varphi_t^*\omega_t)
    &=\lim_{h\to0}\frac{\varphi_{t+h}^*\omega_{t+h}-\varphi_t^*\omega_t}{h} \tag{3.2.3}\\
    &=\lim_{h\to0}\frac{\varphi_{t+h}^*\omega_{t+h}-\varphi_t^*\omega_{t+h}+\varphi_t^*\omega_{t+h}-\varphi_t^*\omega_t}{h} \tag{3.2.4}\\
    &=\lim_{h\to0}\frac{\varphi_{t+h}^*\omega_{t+h}-\varphi_t^*\omega_{t+h}}{h}
      +\lim_{h\to0}\frac{\varphi_t^*\omega_{t+h}-\varphi_t^*\omega_t}{h} \tag{3.2.5}\\
    &=\lim_{h\to0}\frac{\varphi_{t+h}^*\omega_t-\varphi_t^*\omega_t}{h}
      +\varphi_t^*\left(\lim_{h\to0}\frac{\omega_{t+h}-\omega_t}{h}\right) \tag{3.2.6}\\
    &=\varphi_t^*\bigl(\mathcal L_{X_t}\omega_t\bigr)+\varphi_t^*\left(\frac{d}{dt}\omega_t\right) \tag{3.2.7}\\
    &=\varphi_t^*\left(\mathcal L_{X_t}\omega_t+\frac{d}{dt}\omega_t\right) \tag{3.2.8}
  \end{align}
  对任意 $t\in[0,1]$ 为零，即
  \begin{equation}
    \mathcal L_{X_t}\omega_t+\frac{d}{dt}\omega_t=0
    \tag{3.2.9}
  \end{equation}
  成立。这里 $\mathcal L_{X_t}$ 是 $\omega_t$ 关于 $X_t$ 的 Lie 导数，而向量场 $X_t$ 定义为使
  \begin{equation}
    \lim_{h\to0}\frac{\varphi_{t+h}^*\omega_t-\varphi_t^*\omega_t}{h}=\varphi_t^*\bigl(\mathcal L_{X_t}\omega_t\bigr)
    \tag{3.2.10}
  \end{equation}
  成立：
  \begin{equation}
    (X_t)_{\varphi_t(x)}=\left(\frac{d}{ds}\varphi_{t+s}(x)\right)\Big|_{s=0}.
    \tag{3.2.11}
  \end{equation}
  利用 Cartan 同伦公式
  \begin{equation}
    \mathcal L_X=\iota_X\circ d+d\circ\iota_X
    \tag{3.2.12}
  \end{equation}
  以及 $\omega_t$ 为闭形式，可得
  \begin{equation}
    \mathcal L_{X_t}\omega_t=\iota_{X_t}\rd\omega_t+d\iota_{X_t}\omega_t=d\iota_{X_t}\omega_t.
    \tag{3.2.13}
  \end{equation}
  由 $\omega_t$ 所属的 de Rham 上同调类不依赖于 $t$ 的假设，存在某个 1 形式 $u_t$ 使
  \begin{equation}
    \frac{d}{dt}\omega_t=-\rd u_t
    \tag{3.2.14}
  \end{equation}
  % [原文 p.41]
  成立。关于此 $u_t$ 与 $\omega_t$ 对偶的向量场 $X_t$ 由
  \begin{equation}
    u_t=\iota_{X_t}\omega_t
    \tag{3.2.15}
  \end{equation}
  定义。由于 $M$ 紧致，可以对向量场族 $(X_t)_{t\in[0,1]}$ 积分，得到微分同胚族 $(\varphi_t)_{t\in[0,1]}$。该向量场满足 (3.2.9)，故所得的微分同胚族满足 (3.2.2)。
\end{proof}
% [Tutor Note: 此即 Moser 技巧：给定同调的变化路径，用非退化性对偶出向量场，再积分它得到所需的微分同胚；M 紧致保证向量场族可整体积分。]

下面用 Moser 定理证明 Darboux 定理：

\begin{proof}[定理 3.2.2 的证明]
  对任意点 $p\in M$，取 $p$ 的充分小的邻域 $U$，则存在 $\R^n$ 中原点 $0$ 的邻域 $V$，以及把 $0$ 映到 $p$ 的微分同胚 $\psi\colon V\to U$。$M$ 的辛形式 $\omega$ 在 $p$ 处的值 $\omega_p\in\bigwedge^2T_p^*M$ 经 $\psi$ 拉回所得 $\psi^*\omega_p$ 是 $\bigwedge^2T_0^*\R^{2n}$ 的元素；由命题 3.1.1，可适当选取 $\R^{2n}$ 的坐标 $(p_1,\dots,p_n,q_1,\dots,q_n)$，使得
  \[
    \psi^*\omega_p=(\omega_{\R^{2n}})_0
  \]
  成立。其中 $(\omega_{\R^{2n}})_0$ 是由
  \begin{equation}
    \omega_{\R^{2n}}:=\sum_{i=1}^{2n}\rd p_i\wedge \rd q_i
    \tag{3.2.16}
  \end{equation}
  定义的 $\R^{2n}$ 的标准辛形式在 $0\in\R^{2n}$ 处的值。
  % [Tutor Note: (3.2.16) 的求和上限原书作 2n，但 $p_i,q_i$ 只到 $n$，疑为 $n$ 之误，此处照录。]
  这里考虑以 $t\in[0,1]$ 参数化的 $V$ 上 2 次微分形式族
  \begin{equation}
    \omega_t=(1-t)\omega_{\R^{2n}}+t\,\psi^*\omega.
    \tag{3.2.17}
  \end{equation}
  对任意 $t\in[0,1]$，$\omega_t$ 显然是闭形式。又因非退化性是开条件，必要时把 $V$ 取得更小，便可假定对任意 $t\in[0,1]$，$\omega_t$ 非退化。因此，$\omega_t$ 是以 $t\in[0,1]$ 参数化的 $V$ 上辛形式族。取 $V$ 可缩，于是可假定对任意 $t\in[0,1]$，$\omega_t$ 为恰当形式。由 (3.2.14) 与 (3.2.15) 定义向量场族 $(X_t)_{t\in[0,1]}$。若开集 $V'\subset V$ 取得充分小，则可使对 $V'$ 的任意点 $x$，以 $x$ 为初值的 $X_t$ 的积分曲线 $\varphi_t(x)$ 在时刻 1 的值 $\varphi_1(x)$ 仍留在 $V$ 中。由此得到 $V'$ 与 $\varphi_1(V')$ 之间的微分同胚，且由作法
  \begin{equation}
    \varphi_1^*\omega_1=\omega_0
    \tag{3.2.18}
  \end{equation}
  成立。因此，令
  \begin{equation}
    \varphi=\psi\circ\varphi_1\colon V'\to U
    \tag{3.2.19}
  \end{equation}
  并把 $V'$ 重新记作 $V$、把 $\varphi(V')$ 重新记作 $U$，定理 3.2.2 的论断便成立。
  % [Tutor Note: 校勘。原书此处写「$\R^n$ 中原点 $0$ 的邻域 $V$」；但下文 $\psi^*\omega_p$ 落在 $\bigwedge^2T_0^*\R^{2n}$ 中，辛形式也取 $\R^{2n}$ 的标准坐标，故原书此处的 $\R^n$ 应为 $\R^{2n}$。译文照录原文，读者请按 $\R^{2n}$ 理解。]
\end{proof}

% [原文 p.42]
Darboux 定理断言：赋予 $\R^{2n}$ 以标准辛结构所得者，是辛流形任意点邻域的标准形。但已知实际上不仅单个点，任意 Lagrange 子流形的邻域也具有标准形：

\begin{theorem}{3.2.4（Weinstein）}
  设 $(M,\omega)$ 为辛流形，$L$ 为其 Lagrange 子流形。此时存在 $L$ 在 $M$ 中的邻域 $U$ 与 $L$ 在其余切丛 $T^*L$ 中的邻域 $V$，使得 $U$ 与 $V$ 辛同胚。
\end{theorem}
% [Tutor Note: 此处余切丛 $T^*L$ 赋以 4.1 节所述的标准辛结构；$L$ 作为零截面的像自然被视为 $T^*L$ 的 Lagrange 子流形。Weinstein 定理的证明与 Darboux 定理一样借助 Moser 定理，此处省略。]

\section{复流形}

由域的包含关系 $\R\subset\C$，$n$ 维复线性空间 $\C^n$ 可自然视为 $2n$ 维实线性空间 $\R^{2n}$，因此复一般线性群 $\GL(n,\C)$ 自然成为实一般线性群 $\GL(2n,\R)$ 的子群。对 $2n$ 维流形 $M$，把 $TM$ 的结构群约化到 $\GL(n,\C)$ 称为 $M$ 的近复结构（almost complex structure）。给 $M$ 一个近复结构，等价于给 $\Gamma(\End(TM))$ 中一个满足 $J^2=-\id$ 的元素 $J$。给定近复流形 $(M,J_M)$ 与 $(N,J_N)$ 时，满足
\begin{equation}
  \rd\varphi\circ J_M=J_N\circ \rd\varphi
  \tag{3.3.1}
\end{equation}
的映射 $\varphi\colon M\to N$ 称为伪全纯映射（pseudo-holomorphic map）。若近复流形 $(M,J_M)$ 与 $(N,J_N)$ 之间存在伪全纯映射 $\psi\colon M\to N$ 与 $\phi\colon N\to M$ 使 $\psi\circ\phi=\id_N$ 且 $\phi\circ\psi=\id_M$，则称二者作为近复流形同型（isomorphic）。

可积的近复结构称为复结构（complex structure）。换言之，近复流形 $(M,J_M)$ 是复流形，是指对任意 $p\in M$，存在 $p$ 在 $M$ 中的邻域 $U$ 与 $\C^n$ 的开集 $V$，使 $(U,(J_M)|_U)$ 与 $(V,(J_{\C^n})|_V)$ 作为近复流形同型。这里 $J_{\C^n}$ 表示 $\C^n$ 的标准近复结构。复流形之间的伪全纯映射称为全纯映射（holomorphic map）。换言之，把 $\C^n$ 的开集用双全纯映射粘合起来，所得者即为复流形。

对近复流形 $(M,J)$ 与其上的点 $p\in M$，$J_p$ 在 $T_pM$ 上的作用满足 $J_p^2=-\id_{T_pM}$，因此把 $T_pM$ 复化后，可分解为特征值 $\sqrt{-1}$ 的特征空间 $T'_pM$ 与特征值 $-\sqrt{-1}$ 的特征空间 $T''_pM$ 的直和；
% [原文 p.43]
\begin{align}
  T'_pM&=\{X\in T_pM\otimes_\R\C\mid J_pX=\sqrt{-1}X\}, \tag{3.3.2}\\
  T''_pM&=\{X\in T_pM\otimes_\R\C\mid J_pX=-\sqrt{-1}X\}, \tag{3.3.3}\\
  T_pM\otimes_\R\C&=T'_pM\oplus T''_pM. \tag{3.3.4}
\end{align}
各点处的这种直和分解在 $M$ 整体上粘合起来，给出向量丛的直和分解
\begin{equation}
  TM\otimes_\R\C=T'M\oplus T''M.
  \tag{3.3.5}
\end{equation}
$T'M$ 称为 $M$ 的全纯切丛（holomorphic tangent bundle），$T''M$ 称为 $M$ 的反全纯切丛（anti-holomorphic tangent bundle）。此外，用这些向量丛的对偶丛定义
\begin{equation}
  \Omega^{p,q}(M)=\Gamma\left(\bigwedge\nolimits^p(T'M)^*\otimes\bigwedge\nolimits^q(T''M)^*\right),
  \tag{3.3.6}
\end{equation}
则存在直和分解
\begin{equation}
  \Omega^r(M)\otimes\C=\bigoplus_{p+q=r}\Omega^{p,q}(M).
  \tag{3.3.7}
\end{equation}
把附属于此直和分解的投影记为
\begin{equation}
  \Pi^{p,q}\colon\Omega^r(M)\otimes\C\to\Omega^{p,q}(M),
  \tag{3.3.8}
\end{equation}
并定义微分算子 $\partial$ 与 $\overline{\partial}$ 为
\begin{align}
  \partial&:=\Pi^{p+1,q}\circ d\colon\Omega^{p,q}(M)\to\Omega^{p+1,q}(M), \tag{3.3.9}\\
  \overline{\partial}&:=\Pi^{p,q+1}\circ d\colon\Omega^{p,q}(M)\to\Omega^{p,q+1}(M). \tag{3.3.10}
\end{align}
对 $X,Y\in\Gamma(TM)$，令
\begin{equation}
  N(X,Y)=[JX,JY]-[X,Y]-J[X,JY]-J[JX,Y],
  \tag{3.3.11}
\end{equation}
则容易看出对任意 $f\in C^\infty(M)$ 有
\begin{equation}
  N(fX,Y)=N(X,fY)=fN(X,Y),
  \tag{3.3.12}
\end{equation}
故 $N$ 是 $(T^*M)^{\otimes2}\otimes TM$ 的截面。此 $N$ 称为 Nijenhuis 张量（Nijenhuis tensor）。对近复流形 $(M,J)$，下列 3 个条件容易看出彼此等价：
\begin{enumerate}[leftmargin=2.5em,itemsep=2pt,topsep=3pt]
  \item Nijenhuis 张量为零。
  \item 对任意 $X,Y\in\Gamma(T'M)$，$[X,Y]$ 仍属于 $\Gamma(T'M)$。
  \item $\overline{\partial}^2=0$。
\end{enumerate}

% [原文 p.44]
下面的定理是 Darboux 定理在复几何中的对应物。由此定理可知，Nijenhuis 张量相当于辛几何中的 $\rd\omega$。

\begin{theorem}{3.3.1（Newlander--Nirenberg）}
  近复结构 $J$ 可积的充要条件是 Nijenhuis 张量为零。
\end{theorem}
% [Tutor Note: 对比定理 3.2.1：那里「$\omega$ 闭」是可积性条件，这里由 $N=0$ 取代 $\rd\omega=0$，故说 Nijenhuis 张量是 $\rd\omega$ 的复几何类比。]

\takeaway{近复结构可积 $\iff$ Nijenhuis 张量为零（Newlander--Nirenberg）。}

设 $M$ 为复流形，$M$ 的复坐标记为 $(z_i=x_i+\sqrt{-1}y_i)_{i=1}^{n}$，则
\begin{equation}
  J\left(\frac{\partial}{\partial x_i}\right)=\frac{\partial}{\partial y_i},\qquad
  J\left(\frac{\partial}{\partial y_i}\right)=-\frac{\partial}{\partial x_i},
  \tag{3.3.13}
\end{equation}
所以 $M$ 的全纯切丛与反全纯切丛分别由
\begin{equation}
  \frac{\partial}{\partial z_i}=\frac{1}{2}\left(\frac{\partial}{\partial x_i}-\sqrt{-1}\frac{\partial}{\partial y_i}\right),
  \qquad i=1,\dots,n
  \tag{3.3.14}
\end{equation}
与
\begin{equation}
  \frac{\partial}{\partial \bar z_i}=\frac{1}{2}\left(\frac{\partial}{\partial x_i}+\sqrt{-1}\frac{\partial}{\partial y_i}\right),
  \qquad i=1,\dots,n
  \tag{3.3.15}
\end{equation}
张成。它们的对偶基为
\begin{equation}
  \rd z_i=\rd x_i+\sqrt{-1}\,\rd y_i,\qquad
  \rd\bar z_i=\rd x_i-\sqrt{-1}\,\rd y_i.
  \tag{3.3.16}
\end{equation}

\section{驯顺性与相容性}

辛流形 $(M,\omega)$ 上的近复结构 $J$，若对任意 $p\in M$ 与任意非零 $X\in T_pM$ 都有 $\omega_p(X,J_pX)>0$，则称 $J$ 关于 $\omega$ 是驯顺的（tamed by $\omega$）。若 $J$ 关于 $\omega$ 驯顺，则由
\begin{equation}
  g(X,Y)=\frac{1}{2}\bigl(\omega(X,JY)+\omega(Y,JX)\bigr)
  \tag{3.4.1}
\end{equation}
确定的 $g$ 给出 $M$ 的 Riemann 度量。

辛流形 $(M,\omega)$ 上的驯顺近复结构 $J$，若是关于 $g$ 的等距映射，即对任意 $X,Y\in\Gamma(TM)$ 有 $g(JX,JY)=g(X,Y)$，则称 $J$ 与 $\omega$ 相容（compatible with $\omega$）。

\takeaway{驯顺条件 $\omega(X,JX)>0$ 保证由 (3.4.1) 得到的 $g$ 为正定 Riemann 度量。}

\section{Kähler 流形}

复几何与辛几何乍看之下是完全不同类型的几何学，但在某些方面却非常相似。二者都是偶数维，且都可看作实流形在某种意义上「加倍」所得。

% [原文 p.45]
给定实流形 $N$ 时，余切丛 $M=T^*N$ 是辛流形的典型例子，也是分析力学的故乡。其自然的辛结构，在 $N$ 的局部坐标取 $(q_1,\dots,q_n)$、对应的 $T^*N$ 的纤维方向坐标取 $(p_1,\dots,p_n)$ 时，由
\begin{equation}
  \omega=\sum_{i=1}^{n}\rd p_i\wedge \rd q_i
  \tag{3.5.1}
\end{equation}
定义。把此辛形式用坐标 $(p_1,\dots,p_n,q_1,\dots,q_n)$ 表示为交错矩阵，则得
\begin{equation}
  J_{2n}=\begin{pmatrix}0 & I_n\\ -I_n & 0\end{pmatrix}.
  \tag{3.5.2}
\end{equation}
另一方面，切丛 $W=TN$ 上的张量 $J\in C^\infty(W,\End(TW))$，在自然的直和分解 $T(TN)=TN\oplus TN$ 下按与 (3.5.2) 类似的行列
\begin{equation}
  J=\begin{pmatrix}0 & \id_{TN}\\ -\id_{TN} & 0\end{pmatrix}
  \tag{3.5.3}
\end{equation}
定义时，它给出 $W$ 的复结构。

同时具有辛几何与复几何这两种不同几何的，就是 Kähler 流形，其定义如下。

首先注意正交群 $\OO(2n)$、复一般线性群 $\GL(n,\C)$ 与辛群 $\Sp(2n,\R)$ 的公共部分为酉群 $\U(n)$：
\begin{equation}
  \U(n)=\OO(2n)\cap \GL(n,\C)\cap \Sp(2n,\R).
  \tag{3.5.4}
\end{equation}
这可由这些群分别由
\begin{align}
  \U(n)&=\{g\in \GL(n,\C)\mid g^{-1}=g^*\}, \tag{3.5.5}\\
  \OO(2n)&=\{g\in \GL(2n,\R)\mid g^{-1}=g^{\mathsf T}\}, \tag{3.5.6}\\
  \GL(n,\C)&=\{g\in \GL(2n,\R)\mid J_{2n}g=gJ_{2n}\}, \tag{3.5.7}\\
  \Sp(2n,\R)&=\{g\in \GL(2n,\R)\mid g^{\mathsf T}J_{2n}g=J_{2n}\} \tag{3.5.8}
\end{align}
定义，以及从 $n$ 阶复方阵到 $2n$ 阶实方阵的对应中 Hermite 共轭对应于转置，立即推出。因此，具有 $\U(n)$ 结构的流形自动具有 $\OO(2n)$ 结构、$\GL(n,\C)$ 结构与 $\Sp(2n,\R)$ 结构。

\begin{definition}{3.5.1}
  具有 $\U(n)$ 结构、且其附带的 $\GL(n,\C)$ 结构与 $\Sp(2n,\R)$ 结构都可积的流形，称为 Kähler 流形（Kähler manifold）。
\end{definition}

% [原文 p.46]
换言之，Kähler 流形既是复流形又是辛流形，而且这两种结构彼此相容。

在 (3.5.4) 中，右边出现的 3 个群中任意 2 个的公共部分自动包含在剩下的 1 个之中，这一点容易看出。由此，辛结构、复结构与 Riemann 结构这三者中若给定两个，剩下的一个便自动确定。

用刻画它们的张量 $J$、$g$ 与 $\omega$ 来叙述辛结构、近复结构与 Riemann 结构的关系，则为
\begin{equation}
  \omega=g(-,J-).
  \tag{3.5.9}
\end{equation}
实际上，$(J,g,\omega)$ 中给定两个，剩下的一个就由 (3.5.9) 确定。

一般地，近复流形 $(M,J)$ 与其上 Riemann 度量 $g$ 的组，若 $g$ 关于 $J$ 不变，即
\begin{equation}
  g(X,Y)=g(JX,JY),
  \tag{3.5.10}
\end{equation}
则称为近 Hermite 流形（almost Hermitian manifold）。给定近复流形 $(M,J)$ 与其上 Riemann 度量 $g$ 时，把 $g$ 关于 $J$ 作用取平均，便得到近 Hermite 流形：
\begin{equation}
  \tilde g(X,Y):=\frac{1}{2}\bigl(g(X,Y)+g(JX,JY)\bigr).
  \tag{3.5.11}
\end{equation}
近 Hermite 结构等价于未必可积的 $\U(n)$ 结构。近复结构可积的近 Hermite 流形称为 Hermite 流形（Hermitian manifold）。

对近 Hermite 流形 $(M,J,g)$，用 (3.5.9) 定义 2 阶张量 $\omega\in\Gamma\bigl((T^*M)^{\otimes2}\bigr)$，则它成为交错的：
\begin{align}
  \omega(Y,X)&=g(Y,JX) \tag{3.5.12}\\
  &=g(JY,J(JX)) \tag{3.5.13}\\
  &=g(JY,-X) \tag{3.5.14}\\
  &=-g(JY,X) \tag{3.5.15}\\
  &=-g(X,JY) \tag{3.5.16}\\
  &=-\omega(X,Y). \tag{3.5.17}
\end{align}
此 $\omega$ 称为近 Hermite 流形 $(M,J,g)$ 的基本 2 次微分形式（fundamental 2-form）。由 Riemann 度量 $g$ 的非退化性可知基本 2 次微分形式 $\omega$ 非退化，故近 Hermite 流形具有 $\Sp(2n,\R)$ 结构。一般地，此 $\Sp(2n,\R)$ 结构未必可积。由定理 3.2.1，
% [Tutor Note: 此处 (3.5.13) 用了 $g$ 的 $J$ 不变性 $g(Y,JX)=g(JY,J(JX))$，(3.5.16) 用了 $g$ 的对称性。]

% [原文 p.47]
Hermite 流形 $(M,J,g)$ 是 Kähler 流形的充要条件是：由 (3.5.9) 定义的 $(M,J,g)$ 的基本 2 次微分形式为闭形式。此时 $\omega$ 称为 Kähler 形式。

\begin{theorem*}{3.5.2}
  Kähler 流形的复子流形是 Kähler 流形。
\end{theorem*}
% [Tutor Note: 证明要点：把 $M$ 的基本 2 次形式限制到 $N$ 上即得 $N$ 的基本 2 次形式，而限制与外微分可交换，故闭性继承。]

\begin{proof}
  设 $(M,J_M,g_M)$ 为 Kähler 流形，$(N,J_N)\subset(M,J_M)$ 为 $M$ 的复子流形。由定义立即可知，把 $M$ 的基本 2 次微分形式限制到 $N$ 上，所得者与 $N$ 的基本 2 次微分形式一致。由于限制与外微分可交换，$M$ 的基本 2 次微分形式若为闭形式，则 $N$ 的基本 2 次微分形式也为闭形式，定理 3.5.2 得证。
\end{proof}

构造不具有 Kähler 结构的辛流形这一问题，至少可以追溯到 Thurston [128]。由于该文中 Thurston 把此问题归给 Weinstein，[129] 中把此问题称为 Thurston--Weinstein 问题。

\section{辛映射与 Poisson 映射}

对辛流形与 Poisson 流形，可以自然地形式化辛映射与 Poisson 映射的概念：

\begin{definition}{3.6.1}
  对辛流形 $(M,\omega_M)$ 与 $(N,\omega_N)$，满足
  \begin{equation}
    \varphi^*\omega_N=\omega_M
    \tag{3.6.1}
  \end{equation}
  的映射 $\varphi\colon M\to N$ 称为辛映射（symplectic map）。特别地，当 $\varphi$ 为嵌入时，称 $(M,\omega_M)$ 为 $(N,\omega_N)$ 的辛子流形（symplectic submanifold）。
\end{definition}

\begin{definition}{3.6.2}
  Poisson 流形 $(M,\{-,-\}_M)$ 与 $(N,\{-,-\}_N)$ 之间的映射 $\varphi\colon M\to N$，若对任意 $f,g\in C^\infty(N)$ 满足
  \begin{equation}
    \{\varphi^*f,\varphi^*g\}_M=\varphi^*\{f,g\}_N,
    \tag{3.6.2}
  \end{equation}
  则称为 Poisson 映射（Poisson map）。特别地，当 $\varphi$ 为嵌入时，称 $(M,\{-,-\}_M)$ 为 $(N,\{-,-\}_N)$ 的 Poisson 子流形（Poisson submanifold）。
\end{definition}

这些定义虽然自然，却非常强。强到并非没有用处，但如下面两个引理所示，满足这些条件的情形相当有限，这给出了「自然的定义未必是好定义」的典型例子。

% [原文 p.48]
\begin{lemma*}{3.6.3}
  辛流形之间的辛映射
  \[
    \varphi\colon(M,\omega_M)\to(N,\omega_N)
  \]
  存在的必要条件是 $\dim M\le \dim N$。
\end{lemma*}

\begin{proof}
  对任意映射 $\varphi\colon M\to N$、任意点 $p\in M$，以及 $N$ 上任意微分形式 $\omega$，都有
  \begin{equation}
    \rank(\varphi^*\omega)_p\le \rank\omega_{\varphi(p)}
    \tag{3.6.3}
  \end{equation}
  成立。另一方面，$M$ 上 2 次微分形式 $\omega_M$ 非退化的充要条件是：对任意 $p\in M$ 有
  \begin{equation}
    \rank(\omega_M)_p=\dim M.
    \tag{3.6.4}
  \end{equation}
  引理 3.6.3 由这两个事实显然可得。
\end{proof}

\begin{lemma*}{3.6.4}
  设 $(M,\{-,-\}_M)$ 为 Poisson 流形，$N$ 为 $M$ 的子流形。此时，存在 $N$ 上的 Poisson 结构 $\{-,-\}_N$ 使 $(N,\{-,-\}_N)$ 成为 $(M,\{-,-\}_M)$ 的 Poisson 子流形的充要条件是：对 $M$ 上任意函数 $H$ 与 $N$ 上任意点 $p$，由 $H$ 生成的 Hamilton 流
  \begin{equation}
    \varphi\colon\R\to M,\qquad
    \varphi_*\left(\frac{\partial}{\partial t}\right)=(X_H)_{\varphi(t)},\qquad
    \varphi(0)=p
    \tag{3.6.5}
  \end{equation}
  对任意 $t\in\R$ 都满足 $\varphi(t)\in N$。又此时，使 $(N,\{-,-\}_N)$ 成为 $(M,\{-,-\}_M)$ 的 Poisson 子流形的 $N$ 上 Poisson 结构 $\{-,-\}_N$ 是唯一的。
\end{lemma*}

\begin{proof}
  当 $(N,\{-,-\}_N)$ 是 $(M,\{-,-\}_M)$ 的 Poisson 子流形时，对任意 $f,g\in C^\infty(M)$，$\{f,g\}_M$ 只依赖于 $f|_N$ 与 $g|_N$。因此，$M$ 上的向量场 $X_f=\{f,-\}$ 与 $N$ 相切。必要性得证。

  反之，若 $X_f$ 与 $N$ 相切，则 $\{f,g\}_M$ 只依赖于 $g|_N$。交换 $f$ 与 $g$ 的角色，可知 $\{f,g\}_M$ 也只依赖于 $f|_N$。充分性得证。

  唯一性：对任意 $f,g$ 与任意 $p\in N$，取在 $M$ 中 $p$ 的邻域 $U$ 上有定义、且限制到 $U\cap N$ 后与 $f,g$ 一致的函数 $\tilde f,\tilde g$，则必须有
  \begin{equation}
    \{f,g\}_N=\{\tilde f,\tilde g\}_M\big|_N
    \tag{3.6.6}
  \end{equation}
  由此唯一性显然。
\end{proof}

辛流形与 Poisson 流形关系密切，但辛映射与 Poisson 映射的概念却相当不同。例如 $T^*\R^2$ 中的 $T^*\R$ 是辛子流形，却不是 Poisson 子流形。

% [原文 p.49]
\section{Lagrange 子流形}

相空间由在配位空间上添加动量而「加倍」得到，而由动量为零的点构成的相空间的子空间可自然与配位空间等同（即 $M=T^*N$ 的零截面的像可与 $N$ 等同）。此子集具有一个显著性质：$M$ 的辛形式 $\omega$ 限制到它上面为零；并且在具有这种性质的子流形之中，它关于包含关系是极大的。具有这两条性质的 $M$ 的子流形，称为 $M$ 的 Lagrange 子流形：

\begin{definition}{3.7.1}
  辛流形 $(M,\omega)$ 的子流形 $L$ 若满足
  \begin{equation}
    \omega|_L=0\quad\text{且}\quad \dim L=\frac{1}{2}\dim M,
    \tag{3.7.1}
  \end{equation}
  则称为 Lagrange 子流形（Lagrangian submanifold）。
\end{definition}

辛流形的子流形是辛子流形与是 Lagrange 子流形，在某种意义上是两极化的概念。前者要求辛形式的限制非退化，后者要求它为零。前者是开条件，后者是闭条件，因此辛子流形稍微变动后仍是辛子流形，而 Lagrange 子流形稍微变动后一般就不再是 Lagrange 子流形。在这个意义上，Lagrange 子流形是比辛子流形「硬」得多的对象。辛流形的中间维数子流形为 Lagrange 子流形这一条件是一阶非线性偏微分方程，因此对给定的辛流形，一般很难找到它的 Lagrange 子流形。Lagrange 子流形的典型例子，除上述余切丛的截面的像之外，还有可积系统中守恒量的等位面，以及 Kähler 流形的实子流形等。

此外，辛同胚 $\varphi\colon(M_1,\omega_1)\to(M_2,\omega_2)$，通过考虑其图像，可视为直积 $(M_1\times M_2,\pi_1^*\omega_1-\pi_2^*\omega_2)$ 的 Lagrange 子流形的特殊情形。特别地，当辛流形 $M_1$ 与 $M_2$ 分别是 $\R^n$ 的开子流形 $N_1$ 与 $N_2$ 的余切丛，且此 Lagrange 子流形由直积 $N_1\times N_2$ 上的函数 $F$ 的微分 $\rd F$ 的图像给出时，称 $F$ 为正则变换 $\varphi$ 的母函数（generating function）。

记 $M_1$ 的坐标为 $(q,p)$、$M_2$ 的坐标为 $(Q,P)$，则用母函数 $F(q,Q)$，变换前后变量的关系可写作
\begin{equation}
  p_i=\frac{\partial F}{\partial q_i},\qquad
  P_i=-\frac{\partial F}{\partial Q_i}.
  \tag{3.7.2}
\end{equation}
$F$ 是 $q$ 与 $Q$ 的函数，故先解左式即可把 $Q$ 表为 $q$ 与
% [原文 p.50]
$p$ 的函数，再解右式即可把 $P$ 确定为 $q$ 与 $p$ 的函数。右式的符号是为与分析力学中的标准约定一致而设，可通过把 $P_i$ 乘以 $-1$ 等方法吸收，故并非本质。

辛几何中许多重要概念都可用 Lagrange 子流形的语言书写，Weinstein 因此提出「一切都是 Lagrange 子流形」的口号。这就是 Weinstein 的辛几何信条（Weinstein's symplectic creed）[137]。虽说「一切」有些夸大，但 Lagrange 子流形的研究是辛几何的中心问题，这一点毫无疑问。
% [Tutor Note: 此信条出自 Weinstein 1981 年的论文《Symplectic geometry》，意谓辛几何的对象与构造都可用 Lagrange 子流形重新表述。]

\takeaway{Lagrange 子流形是 $\omega$ 的限制为零且维数为一半的子流形，是辛几何的中心对象。}

\chapterreview{%
本章以 Klein 的埃尔朗根纲领为引子，把几何学理解为空间与其变换群的组，并用 Cartan 的 G 结构把这一观点推广到流形：辛结构即结构群约化到 $\Sp(2n,\R)$ 的 G 结构，而复结构则是约化到 $\GL(n,\C)$ 的结构。Darboux 定理（及其 Moser 证明）表明辛结构局部都是标准形，因此辛几何研究的是辛同胚不变的整体性质；Weinstein 定理进一步说明 Lagrange 子流形的邻域也有标准形。复几何一侧，Newlander--Nirenberg 定理把近复结构的可积性归结为 Nijenhuis 张量为零，与辛几何中 $\rd\omega=0$ 的地位相对应。Kähler 流形同时具有复结构、辛结构与相容的 Riemann 度量，其基本 2 次形式闭；而一般的辛流形未必容许 Kähler 结构。最后引入的辛映射与 Poisson 映射虽然自然，却是很强的条件，这与 Lagrange 子流形作为「硬」对象的地位一起，构成后续各章讨论辛商与可积系统的出发点。%
}
